\chapter*{\hfill{\centering  ЗАКЛЮЧЕНИЕ}\hfill}
\addcontentsline{toc}{chapter}{ЗАКЛЮЧЕНИЕ}

В ходе исследований получено, что максимальное число сравнений при использовании реализации бинарного поиска равно 10, а при использовании реализации линейного поиска~---~1024. Среднее количество сравнений для нахождения элемента в массиве алгоритмом бинарного поиска~---~8, линейного поиска~---~393. Алгоритм бинарного поиска лучше линейного поиска при большем размере исходного массива, поскольку требует в среднем меньше сравнений. При использовании бинарного поиска стоит учитывать необходимость упорядоченности элементов массива.

Поставленная цель была достигнута.
Для достижения поставленной цели были выполнены все поставленные задачи:
\begin{itemize}[label=---]
	\item описаны линейный и бинарный алгоритмы поиска;
	\item написана программа, реализующая алгоритм поиска полным перебором и бинарный алгоритм поиска;
	\item проверена работоспособность программы, проведено модульное тестирование;
	\item замерено число сравнений элементов в упомянутых выше алгоритмах при разном размере массива и положении искомого элемента;
	\item проведён анализ полученных результатов.
\end{itemize}
